% BEGIN VOL07-EXPANSION VII11
\section{Supplementary worked examples}
\label{sec:vii11-pedagogy}

\begin{example}[The fundamental theorem of calculus is Stokes in dimension one]\label{ex:vii11-ped-01}
For an oriented interval \([a,b]\) and a smooth function \(f\),
\[
\int_{[a,b]}df=f(b)-f(a)=\int_{\partial[a,b]}f.
\]
The boundary is the signed zero-chain \(\{b\}-\{a\}\).  Stokes' theorem therefore contains the fundamental theorem of calculus as its first case.
\end{example}

\begin{example}[Green's theorem from Stokes]\label{ex:vii11-ped-02}
For \(\alpha=P\,dx+Q\,dy\) on an oriented planar region \(D\),
\[
d\alpha=(Q_x-P_y)\,dx\wedge dy.
\]
Stokes gives
\[
\int_{\partial D}P\,dx+Q\,dy=\int_D(Q_x-P_y)\,dx\,dy,
\]
which is Green's theorem with the positive boundary orientation.
\end{example}

\begin{example}[Exact forms have zero integral on closed manifolds]\label{ex:vii11-ped-03}
If \(M\) is compact, oriented, and has no boundary, then for every \((n-1)\)-form \(\eta\),
\[
\int_M d\eta=\int_{\partial M}\eta=0.
\]
This converts a geometric boundary statement into an immediate obstruction to exactness.
\end{example}

\section{Graded supplementary exercises}
The following exercises add a deliberate progression from concrete calculation to structural proof, hypothesis testing, application, and synthesis.

\subsection*{Standard computations and constructions}

\begin{exercise}[Interval Stokes]\label{exr:vii11-ped-01}
Apply Stokes to \(f(x)=x^2\) on \([0,2]\).
\end{exercise}
\begin{hint}
Integrate \(df=2x\,dx\).
\end{hint}
\begin{solution}
\(\int_0^2 2x\,dx=4=f(2)-f(0)\).
\end{solution}

\begin{exercise}[Rectangle circulation]\label{exr:vii11-ped-02}
For \(\alpha=x\,dy\), compute \(\int_{\partial[0,1]^2}\alpha\) using Stokes.
\end{exercise}
\begin{hint}
Compute \(d\alpha\).
\end{hint}
\begin{solution}
\(d\alpha=dx\wedge dy\), so the boundary integral equals the area \(1\).
\end{solution}

\begin{exercise}[Exact one-form on a circle]\label{exr:vii11-ped-03}
Compute \(\int_{S^1} d(x|_{S^1})\).
\end{exercise}
\begin{hint}
The circle has no boundary.
\end{hint}
\begin{solution}
The integral is \(0\) by Stokes, since the form is exact on the closed one-manifold.
\end{solution}

\begin{exercise}[Annulus boundary signs]\label{exr:vii11-ped-04}
State the induced orientations on the outer and inner boundary circles of an oriented planar annulus.
\end{exercise}
\begin{hint}
Use outward-normal-first convention.
\end{hint}
\begin{solution}
The outer circle is counterclockwise and the inner circle is clockwise.
\end{solution}

\begin{exercise}[Disk angular computation]\label{exr:vii11-ped-05}
For \(\alpha=x\,dy-y\,dx\), compute \(\int_{\partial D}\alpha\) on the unit disc.
\end{exercise}
\begin{hint}
Compute \(d\alpha\).
\end{hint}
\begin{solution}
\(d\alpha=2\,dx\wedge dy\), so the integral is \(2\pi\).
\end{solution}

\subsection*{Proofs}

\begin{exercise}[Exact integral vanishes]\label{exr:vii11-ped-06}
Prove that the integral of an exact top-degree form over a compact oriented manifold without boundary is zero.
\end{exercise}
\begin{hint}
Write the form as \(d\eta\) and apply Stokes.
\end{hint}
\begin{solution}
\(\int_M d\eta=\int_{\partial M}\eta=0\).
\end{solution}

\begin{exercise}[Boundary of boundary]\label{exr:vii11-ped-07}
Explain why Stokes is compatible with the principle \(\partial^2=0\).
\end{exercise}
\begin{hint}
Apply Stokes twice to a form of degree \(n-2\).
\end{hint}
\begin{solution}
\(\int_{\partial(\partial M)}\eta=\int_{\partial M}d\eta=\int_Md^2\eta=0\), matching the geometric fact that an oriented boundary has no boundary.
\end{solution}

\begin{exercise}[Orientation reversal]\label{exr:vii11-ped-08}
Show Stokes remains valid after reversing the orientation of \(M\).
\end{exercise}
\begin{hint}
Track both the interior and induced boundary orientations.
\end{hint}
\begin{solution}
Reversing \(M\) negates the left integral and also reverses the induced boundary orientation, negating the right integral.
\end{solution}

\begin{exercise}[Closed form period obstruction]\label{exr:vii11-ped-09}
Prove that if a one-form has nonzero integral around a closed curve, then it is not exact.
\end{exercise}
\begin{hint}
Assume \(\omega=df\) and apply Stokes to the closed one-manifold.
\end{hint}
\begin{solution}
An exact one-form integrates to zero on every closed curve, so a nonzero period contradicts exactness.
\end{solution}

\subsection*{Counterexamples and hypothesis tests}

\begin{exercise}[Boundary orientation is arbitrary]\label{exr:vii11-ped-10}
Test this claim.
\end{exercise}
\begin{hint}
Use the interval or disc.
\end{hint}
\begin{solution}
False.  Stokes fixes the induced boundary orientation; changing it alone changes the sign of the boundary integral and breaks the formula.
\end{solution}

\begin{exercise}[Stokes needs no regularity or compactness control]\label{exr:vii11-ped-11}
Test the unrestricted statement on noncompact domains.
\end{exercise}
\begin{hint}
Recall the usual compact-support or integrability hypotheses.
\end{hint}
\begin{solution}
False as stated.  One needs hypotheses ensuring both integrals are defined, such as compactness or compact support in standard formulations.
\end{solution}

\begin{exercise}[Closed implies exact by Stokes]\label{exr:vii11-ped-12}
Test on the punctured plane.
\end{exercise}
\begin{hint}
Use the angular form with nonzero period.
\end{hint}
\begin{solution}
False.  Stokes detects exactness obstructions but does not make every closed form exact on topologically nontrivial spaces.
\end{solution}

\subsection*{Applications and investigations}

\begin{exercise}[Green's theorem]\label{exr:vii11-ped-13}
Derive Green's theorem from Stokes for \(\alpha=Pdx+Qdy\).
\end{exercise}
\begin{hint}
Compute \(d\alpha\).
\end{hint}
\begin{solution}
Since \(d\alpha=(Q_x-P_y)dx\wedge dy\), Stokes gives the standard circulation formula.
\end{solution}

\begin{exercise}[Flux form]\label{exr:vii11-ped-14}
Explain how a divergence theorem can be viewed as Stokes after converting a vector field into an \((n-1)\)-form using a volume form.
\end{exercise}
\begin{hint}
Contract the volume form with the vector field.
\end{hint}
\begin{solution}
If \(\eta=\iota_X\Omega\), then \(d\eta=(\operatorname{div}X)\Omega\); Stokes turns the volume integral of divergence into boundary flux.
\end{solution}

\subsection*{Challenge problems}

\begin{exercise}[Punctured-plane period]\label{exr:vii11-ped-15}
Use Stokes to explain why the angular form on \(\mathbb R^2\setminus\{0\}\) cannot be exact.
\end{exercise}
\begin{hint}
Integrate it on the unit circle.
\end{hint}
\begin{solution}
Its period is \(2\pi\).  Were it exact, the integral on the closed circle would vanish, contradiction.
\end{solution}

\begin{exercise}[Two homologous boundary curves]\label{exr:vii11-ped-16}
Suppose an oriented surface region has boundary \(C_2-C_1\) and \(d\omega=0\).  Show \(\int_{C_2}\omega=\int_{C_1}\omega\).
\end{exercise}
\begin{hint}
Apply Stokes to the region.
\end{hint}
\begin{solution}
\(0=\int_Md\omega=\int_{\partial M}\omega=\int_{C_2}\omega-\int_{C_1}\omega\).
\end{solution}

% END VOL07-EXPANSION VII11
